\documentclass[11pt,a4paper]{article}
\usepackage{shared/luxcover}
\usepackage[utf8]{inputenc}
\usepackage[T1]{fontenc}
\usepackage{amsmath,amssymb,amsfonts,amsthm}
\usepackage{graphicx}
\usepackage{hyperref}
\usepackage{algorithm}
\usepackage{algpseudocode}
\usepackage{booktabs}
\usepackage{listings}
\input{shared/lstlang}
\usepackage{xcolor}
\usepackage{enumitem}
\usepackage{float}

\newtheorem{theorem}{Theorem}
\newtheorem{lemma}[theorem]{Lemma}
\newtheorem{definition}{Definition}

\lstset{
    basicstyle=\small\ttfamily,
    keywordstyle=\color{blue}\bfseries,
    commentstyle=\color{gray},
    stringstyle=\color{red},
    breaklines=true,
    frame=single,
    numbers=left,
    numberstyle=\tiny\color{gray}
}

\hypersetup{
    colorlinks=true,
    linkcolor=blue,
    citecolor=blue,
    urlcolor=blue
}

\title{SR25519 Precompile:\\Schnorrkel Signatures for Substrate Compatibility}
\author{
Zach Kelling\thanks{Corresponding author: zach@lux.network}\\
\textit{Lux Industries, Inc.}\\
\texttt{research@lux.network}
}
\date{April 2026}

\begin{document}

\luxcoverpage

\begin{abstract}
We specify the SR25519 (Schnorrkel) signature verification precompile on the
Lux EVM. SR25519 is the Schnorr-based signature scheme on Ristretto25519
used by Polkadot, Kusama, and all Substrate-based chains for account signing,
session keys, and VRF (Verifiable Random Function) output. Unlike Ed25519,
SR25519 uses the Ristretto group encoding (eliminating cofactor complexities),
produces non-malleable signatures, and supports VRF output verification as a
built-in capability. The precompile provides signature verification and VRF
output verification. Gas cost: 2\,500 for signature verify, 4\,000 for VRF verify.
\end{abstract}

\section{Introduction}

SR25519~\cite{schnorrkel} (Schnorrkel) is the signature scheme of the Substrate
ecosystem. It is used by:
\begin{itemize}
  \item Polkadot relay chain and parachains (validator and session keys).
  \item Kusama (canary network).
  \item All Substrate-based chains (default account signing).
\end{itemize}

SR25519 differs from Ed25519 in three critical ways:
\begin{enumerate}[label=(\roman*)]
  \item \textbf{Ristretto encoding}: Uses the Ristretto255 group~\cite{ristretto}
    (a prime-order group constructed from Curve25519), eliminating cofactor
    complications and signature malleability.
  \item \textbf{Schnorr signature}: Uses a Schnorr construction with
    Merlin transcript-based challenges (domain separation, streaming).
  \item \textbf{VRF support}: Built-in verifiable random function (VRF) via
    the same key pair, used for Polkadot's BABE block production.
\end{enumerate}

\section{Ristretto255}

\begin{definition}[Ristretto255]
Ristretto255 is a prime-order group of order $\ell = 2^{252} + 27742317777372353535851937790883648493$
constructed from Curve25519 via a quotient encoding. Every Ristretto point has
a unique canonical encoding (32 bytes). The cofactor is eliminated by the encoding,
not by multiplication.
\end{definition}

The Ristretto encoding solves the ``cofactor problem'' that plagues Ed25519:
there is exactly one valid encoding for each group element, and decoding
either produces a valid group element or fails. No special-case handling for
low-order points, identity, or non-canonical encodings is needed.

\section{Precompile Specification}

Address: \texttt{0x0200\ldots2E}.

\begin{table}[H]
\centering
\begin{tabular}{llr}
\toprule
\textbf{Selector} & \textbf{Operation} & \textbf{Gas} \\
\midrule
\texttt{0x01} & Signature Verify & 2\,500 \\
\texttt{0x02} & VRF Verify & 4\,000 \\
\texttt{0x03} & Batch Verify & $2\,500 + 1\,800(n-1)$ \\
\bottomrule
\end{tabular}
\end{table}

\subsection{Signature Verification}

\textbf{Input}: Selector (1), public key (32 bytes, Ristretto encoding),
signature (64 bytes: 32-byte $R$, 32-byte $s$), signing context length (2 bytes),
signing context, message length (4 bytes), message.

\textbf{Output}: 32 bytes, value 1 if valid, 0 otherwise.

\textbf{Verification}:
\begin{enumerate}
  \item Decode $R$ as Ristretto point (fail if non-canonical).
  \item Decode public key $A$ as Ristretto point.
  \item Construct Merlin transcript with signing context and message.
  \item Derive challenge $k$ from transcript.
  \item Check $[s]B = R + [k]A$.
\end{enumerate}

The Merlin transcript~\cite{merlin} provides domain separation: the signing
context (a byte string, typically ``substrate'') is committed before the message,
preventing cross-protocol signature replay.

\subsection{VRF Verification}

\textbf{Input}: Selector (1), public key (32 bytes), VRF output (32 bytes),
VRF proof (64 bytes), transcript data (variable).

\textbf{Output}: 32 bytes, value 1 if the VRF output is valid for the given
input and public key.

The VRF~\cite{vrf-schnorrkel} produces a deterministic output $H = [\text{sk}] \cdot \text{HashToPoint}(\text{input})$
that is unpredictable without the secret key but verifiable with the public key.

\section{Security Analysis}

\subsection{Non-Malleability}

SR25519 signatures are non-malleable by construction:
\begin{itemize}
  \item The Ristretto encoding ensures unique point representation.
  \item The scalar $s$ is checked to be in $[0, \ell)$.
  \item The Merlin transcript commits all inputs in order.
\end{itemize}

\subsection{Signing Context}

The mandatory signing context prevents cross-protocol attacks. A signature
produced with context ``substrate'' is invalid under context ``polkadot''.
The precompile requires the signing context as input and rejects empty contexts.

\subsection{VRF Security}

The VRF output is indistinguishable from random under the Decisional
Diffie-Hellman (DDH) assumption on Ristretto255. The proof ensures that the
output was computed correctly by the claimed public key holder.

\section{Gas Cost Justification}

SR25519 verification requires 2 Ristretto scalar multiplications (same as Ed25519)
plus Merlin transcript operations (SHA-3 based, negligible cost). The 25\% premium
over Ed25519 (2\,500 vs 2\,000) reflects the Ristretto decoding overhead and
Merlin transcript construction.

\begin{table}[H]
\centering
\begin{tabular}{lrr}
\toprule
\textbf{Operation} & \textbf{Time ($\mu$s)} & \textbf{Gas} \\
\midrule
Signature Verify & 40 & 2\,500 \\
VRF Verify & 62 & 4\,000 \\
Batch Verify (16) & 420 & 29\,700 \\
\bottomrule
\end{tabular}
\end{table}

\section{Conclusion}

The SR25519 precompile enables native verification of Substrate ecosystem
signatures on the Lux EVM. The Ristretto encoding eliminates cofactor-related
vulnerabilities, and the built-in VRF support enables on-chain verification
of Polkadot BABE block production proofs.

\begin{thebibliography}{9}
\bibitem{schnorrkel}
  J.~Burdges, ``Schnorrkel: Schnorr VRFs and signatures on the Ristretto group,''
  W3F Research, 2019. \url{https://github.com/w3f/schnorrkel}

\bibitem{ristretto}
  H.~de Valence et al., ``Ristretto: prime order elliptic curve groups with
  non-malleable encodings,'' 2018. \url{https://ristretto.group}

\bibitem{merlin}
  H.~de Valence, ``Merlin: composable proof transcripts for public-coin
  arguments of knowledge,'' 2018. \url{https://merlin.cool}

\bibitem{vrf-schnorrkel}
  J.~Burdges, ``A Schnorr VRF with an extra Schnorr DLEQ proof,''
  W3F Research, 2020.
\end{thebibliography}

\end{document}
